\section{Experimental progression}

\headingindent
The initial CIFAR experiments explored many combinations of teacher
architecture, student model, feature layer, perturbation, training objective,
inference protocol, and OOD Score. For brevity, this report retains only the
best-performing variant within each method family and ID dataset.
Feature-denoising experiments are reported separately in Section~5.

AUROC is the main ranking metric used in OOD detection. Under the
ID-positive convention used here, it measures how often a randomly selected ID
sample receives a higher score than a randomly selected OOD sample. It is
threshold-independent. FPR@95 instead describes a particular operating point:
the threshold is chosen to correctly accept 95\% of ID samples, and the metric
reports the proportion of OOD samples that incorrectly pass that threshold and
are accepted as ID. Thus, higher AUROC and lower FPR@95 are better.

% Each table follows the grouped layout used in the full results report. The
% opposite CIFAR dataset is Near-OOD, while MNIST and SVHN form the Far-OOD
% group. The Far-OOD average is the mean over MNIST and SVHN, and the overall
% macro is the unweighted mean over all three OOD datasets. Every metric cell
% reports \textbf{AUROC / FPR@95} under the ID-positive convention. Teacher
% baselines always use maximum softmax probability (MSP), never energy.

\clearpage
\begin{landscape}

\subsection{CIFAR-10 as ID}

\headingindent
Table~\ref{tab:progression-cifar10} reports the selected configuration for each
method family. 
% Dropout and multi-layer clipping were the strongest
% student-based variants by the compromise criterion. Their gains came from
% different inference protocols: the selected dropout model averaged perturbed
% draws, whereas the selected clipping model used clipped inference after
% training against a clean teacher target.

\begin{table}[H]
\centering
\small
\setlength{\tabcolsep}{3.2pt}
\renewcommand{\arraystretch}{1.12}
\begin{tabularx}{\linewidth}{@{}L{2.25cm}YL{2.15cm}ccccc@{}}
\toprule
\multicolumn{3}{c}{} &
\multicolumn{1}{c}{\cellcolor{neargreen}\textbf{Near-OOD}} &
\multicolumn{3}{c}{\cellcolor{farorange}\textbf{Far-OOD}} &
\multicolumn{1}{c}{\cellcolor{macroblue}\textbf{Overall}} \\
\cmidrule(lr){4-4}\cmidrule(lr){5-7}\cmidrule(lr){8-8}
Method family & Best-performing variant & Selected OOD Score &
\cellcolor{neargreen}CIFAR-100 &
\cellcolor{farorange}MNIST & \cellcolor{farorange}SVHN &
\cellcolor{faraverage}Average & \cellcolor{macroblue}Macro \\
\midrule
Teacher baseline & ResNet-18; no student & Teacher MSP &
\nearcell{\metric{0.879}{0.624}} &
\farcell{\metric{0.918}{0.556}} & \farcell{\metric{0.904}{0.598}} &
\faravg{\metric{0.911}{0.577}} & \balanced{\metric{0.900}{0.593}} \\
Teacher baseline & ResNet-50; no student & Teacher MSP &
\nearcell{\metric{0.873}{0.628}} &
\farcell{\metric{0.897}{0.622}} & \farcell{\metric{0.891}{0.680}} &
\faravg{\metric{0.894}{0.651}} & \balanced{\metric{0.887}{0.643}} \\
\addlinespace
Clean distillation & ResNet-18 layer 4; 100-tree random forest; output
regression & Student energy &
\nearcell{\metric{0.870}{0.550}} &
\farcell{\metric{0.920}{0.490}} & \farcell{\metric{0.900}{0.550}} &
\faravg{\metric{0.910}{0.520}} & \balanced{\metric{0.895}{0.531}} \\
MC dropout & ResNet-18 layer 4; channel dropout $p=0.8$; logit MSE; dropout
inference & Student energy &
\nearcell{\metric{0.885}{0.498}} &
\farcell{\metric{0.958}{0.252}} & \farcell{\metric{0.930}{0.380}} &
\faravg{\metric{0.944}{0.316}} & \balanced{\metric{0.924}{0.376}} \\
Multi-layer clipping & ResNet-18 layers 3--4; channel/flatten; clean target;
KL; clipped inference & Student energy &
\nearcell{\metric{0.859}{0.581}} &
\farcell{\metric{0.998}{0.001}} & \farcell{\metric{0.897}{0.482}} &
\faravg{\metric{0.948}{0.242}} & \balanced{\metric{0.918}{0.355}} \\
Affine pixel perturbation & ResNet-50 layer 4; clean target; logit MSE; clean
inference & Student energy &
\nearcell{\metric{0.871}{0.552}} &
\farcell{\metric{0.906}{0.528}} & \farcell{\metric{0.960}{0.247}} &
\faravg{\metric{0.933}{0.388}} & \balanced{\metric{0.912}{0.442}} \\
Subspace-student ensemble & 16 ResNet-18 linear students; ordered disjoint
channel-GAP subspaces; KL & Predictive entropy &
\nearcell{\metric{0.884}{0.574}} &
\farcell{\metric{0.930}{0.468}} & \farcell{\metric{0.905}{0.554}} &
\faravg{\metric{0.918}{0.511}} & \balanced{\metric{0.906}{0.532}} \\
\bottomrule
\end{tabularx}
\caption{Best-performing variant in each method family with CIFAR-10 as ID.
Cells report AUROC / FPR@95 on the initial CIFAR protocol.}
\label{tab:progression-cifar10}
\end{table}

\clearpage
\vspace*{0.8cm}
\subsection{CIFAR-100 as ID}

\headingindent
Table~\ref{tab:progression-cifar100} repeats the selection independently with
CIFAR-100 as ID.
% This setting was consistently harder. Affine pixel
% perturbation and multi-layer clipping gave the strongest student-based
% compromises, while dropout and the subspace-student ensemble were less stable
% than on CIFAR-10.

\begin{table}[H]
\centering
\small
\setlength{\tabcolsep}{3.2pt}
\renewcommand{\arraystretch}{1.12}
\begin{tabularx}{\linewidth}{@{}L{2.25cm}YL{2.15cm}ccccc@{}}
\toprule
\multicolumn{3}{c}{} &
\multicolumn{1}{c}{\cellcolor{neargreen}\textbf{Near-OOD}} &
\multicolumn{3}{c}{\cellcolor{farorange}\textbf{Far-OOD}} &
\multicolumn{1}{c}{\cellcolor{macroblue}\textbf{Overall}} \\
\cmidrule(lr){4-4}\cmidrule(lr){5-7}\cmidrule(lr){8-8}
Method family & Best-performing variant & Selected OOD Score &
\cellcolor{neargreen}CIFAR-10 &
\cellcolor{farorange}MNIST & \cellcolor{farorange}SVHN &
\cellcolor{faraverage}Average & \cellcolor{macroblue}Macro \\
\midrule
Teacher baseline & ResNet-18; no student & Teacher MSP &
\nearcell{\metric{0.792}{0.794}} &
\farcell{\metric{0.701}{0.956}} & \farcell{\metric{0.806}{0.812}} &
\faravg{\metric{0.754}{0.884}} & \balanced{\metric{0.767}{0.854}} \\
Teacher baseline & ResNet-50; no student & Teacher MSP &
\nearcell{\metric{0.792}{0.781}} &
\farcell{\metric{0.811}{0.823}} & \farcell{\metric{0.758}{0.793}} &
\faravg{\metric{0.785}{0.808}} & \balanced{\metric{0.787}{0.799}} \\
\addlinespace
Clean distillation & ResNet-18 layer 4; 100-tree random forest; output
regression & Student energy &
\nearcell{\metric{0.790}{0.810}} &
\farcell{\metric{0.710}{0.940}} & \farcell{\metric{0.840}{0.630}} &
\faravg{\metric{0.775}{0.785}} & \balanced{\metric{0.780}{0.793}} \\
MC dropout & ResNet-18 layer 4; spatial dropout $p=0.8$; KL; clean inference &
KL mismatch &
\nearcell{\metric{0.768}{0.881}} &
\farcell{\metric{0.772}{0.821}} & \farcell{\metric{0.812}{0.854}} &
\faravg{\metric{0.792}{0.838}} & \balanced{\metric{0.784}{0.852}} \\
Multi-layer clipping & ResNet-18 layers 3--4; channel/flatten; clean target;
KL; clean inference & Student energy &
\nearcell{\metric{0.751}{0.842}} &
\farcell{\metric{0.797}{0.879}} & \farcell{\metric{0.940}{0.340}} &
\faravg{\metric{0.869}{0.610}} & \balanced{\metric{0.829}{0.687}} \\
Affine pixel perturbation & ResNet-50 layer 4; clean target; logit MSE; clean
inference & Student energy &
\nearcell{\metric{0.779}{0.797}} &
\farcell{\metric{0.875}{0.698}} & \farcell{\metric{0.925}{0.431}} &
\faravg{\metric{0.900}{0.565}} & \balanced{\metric{0.860}{0.642}} \\
Subspace-student ensemble & 16 ResNet-18 linear students; ordered disjoint
flattened-channel subspaces; logit MSE & Predictive entropy &
\nearcell{\metric{0.779}{0.818}} &
\farcell{\metric{0.760}{0.947}} & \farcell{\metric{0.874}{0.688}} &
\faravg{\metric{0.817}{0.818}} & \balanced{\metric{0.804}{0.818}} \\
\bottomrule
\end{tabularx}
\caption{Best-performing variant in each method family with CIFAR-100 as ID.
Cells report AUROC / FPR@95 on the initial CIFAR protocol.}
\label{tab:progression-cifar100}
\end{table}

Across both ID datasets, no student family dominated every setting. The
selected variant changed with the ID dataset, and strong AUROC did not always
coincide with low FPR@95. This instability motivated the later emphasis on
feature reconstruction pipeline.
% and on fixed OpenOOD evaluation.

\end{landscape}
\clearpage
